\documentclass[../../main2.tex]{subfiles}

\begin{document}

\chapter*{致谢}
\addcontentsline{toc}{chapter}{致谢}

\begin{motivation}
	一本书的诞生，本身就是一次「微观动机汇聚成宏观结构」——正好用这本书自己的话，记下这条通道上的人。\\
	\textbf{核心动机}：感谢，也示范——谢意也要按「谁的贡献拿不掉」来写。
\end{motivation}

按无你检验的规矩，逐个做「拿掉会怎样」：

\begin{itemize}
	\item \textbf{拿掉最早那批读者}：旧稿《社会动力学与公理化体系》的批评者们——尤其是指出「附录断裂」「四章像合订本」「只会解释不会预测」的那几位。$C$ 很大：没有这三刀，就没有这一版的前向结构。
	\item \textbf{拿掉「太不自然」那句否决}：v1 架构里那层「演化状态空间」会一直躺到印厂。否决它的人坚持：读者的问题比哲学家的问题优先。这条原则现在写进了每一节的开头。
	\item \textbf{拿掉法律史}：本书的核心工具不是发明，是抄法庭的作业。所有被「必须给数字」逼了几千年的法律人——你们才是这套方法的原作者。
	\item \textbf{拿掉家人的容忍}：一个盯着「菜价为什么涨」发呆的人，需要有人把菜买回来。$C$ 难估，但不为零。
	\item \textbf{拿掉你}：读到这里的人。一本书没人读，只是一叠纸；有人读、有人拆、有人转给下一个人，它才是书。
\end{itemize}

交互项记一笔：以上任何单独一项都未必能成书——合起来才成。这块「合谋」的功劳，按第 9 章的规矩，单列，均摊给所有被我漏掉的名字。

写作上，本书的写法守则来自我自己的对话语料（短句、第一人称落在决策上、禁空话）；结构上，六段式链条（动机 $\to$ 因果 $\to$ 法庭 $\to$ 齿轮 $\to$ 循环）是和早期评审反复吵出来的。凡书中有错，账记我头上——按贡献算，这口锅我的 $C$ 最大。

\paragraph*{逻辑衔接}
	谢完了。附录的公式、正文的链条、失败的清单，都摆在那里——接下来的事，是你的反事实推演。

\begin{thinkbox}
\begin{enumerate}
	\item 用「拿掉会怎样」给你自己的关键帮助者排个序：谁的边际贡献最大？想一个人，把感谢说出口——感谢也是模因，转出去才算数。
	\item \textbf{反事实}：如果当年那句「太不自然」没有说出口，这本书会是什么形状？你会更喜欢那个版本，还是这一个？——顺手想想：你手上有哪个「太不自然」的方案，其实该被砍掉？
\end{enumerate}
\end{thinkbox}

\end{document}
